\chapter{Thực nghiệm, kết quả và phân tích}
\label{chap:results}

Chương này đối chiếu các mô hình từ MRI một thời điểm, MRI dọc theo thời gian đến các cách hợp nhất thông tin, sau đó dùng ablation để xác định những thành phần cần giữ trong kiến trúc cuối. Kết quả validation được dùng để chọn mô hình; kết quả test hậu nghiệm giúp mô tả hiệu năng của checkpoint đã chọn. RC-Free Pairwise Fusion được giữ làm kiến trúc chính vì có validation cao nhất ở phép so sánh cùng seed và chỉ cần hai cặp tương tác.

\section{Giao thức thực nghiệm và cách đọc kết quả}
\label{sec:experimental-protocol}
Cohort đóng băng gồm 359 người và 1.077 lượt khám, đúng ba visit cho mỗi người. Split theo người bệnh gồm 251 người train, 54 người validation và 54 người test, với 74, 16 và 16 biến cố tương ứng. Chuẩn hóa và PCA chỉ được fit trên train, rồi áp dụng các artifact đã đóng băng cho validation và test. Đầu ra là raw risk; giá trị lớn hơn biểu thị nguy cơ chuyển đổi sớm hơn.

Các bảng một seed sử dụng seed 20260727 và các checkpoint được chọn theo validation trước khi khóa cấu hình để đánh giá test, theo báo cáo đánh giá từ máy chủ. Các bảng ghi epoch của checkpoint được chọn, không phải tổng số epoch huấn luyện. Chỉ số là Harrell-style C-index: người có biến cố được so với người có thời gian theo dõi dài hơn, risk bằng nhau nhận 0,5 điểm và thời gian bằng nhau không tạo cặp so sánh \cite{harrell1996multivariable}. Validation có 480 cặp hợp lệ, còn test có 528 cặp. C-index càng cao càng tốt, được ký hiệu bằng $\uparrow$; ô in đậm là giá trị cao nhất của từng cột Val hoặc Test trong bảng.

Tập test đã từng được truy cập trong quá trình nghiên cứu, vì vậy toàn bộ cột Test trong chương này là \emph{test hậu nghiệm}. Chúng không được dùng để chọn kiến trúc và không được xem là đánh giá độc lập cuối cùng. Benchmark ba seed ở Mục~\ref{sec:standardized-benchmark} chỉ báo cáo validation; trung bình ba seed và một điểm test của seed 20260727 là hai loại bằng chứng khác nhau. Báo cáo test mới được lưu cùng nguồn số liệu, còn các dự đoán từng người và checkpoint trên máy chủ chưa được kiểm tra lại tại môi trường biên soạn.

Với họ MedicalNet Pairwise, CP Pairwise và RC-Free Pairwise, giao thức chính dùng AdamW, learning rate $10^{-4}$, weight decay $10^{-4}$, FP16 AMP, clipping norm 1 và tối đa 100 epoch \cite{loshchilov2019adamw}. ReduceLROnPlateau theo Cox loss toàn validation; early stopping theo C-index dùng warmup 10 epoch, patience 20 và min-delta 0,002. Cox loss xử lý biến cố đồng thời bằng Breslow \cite{cox1972regression,breslow1974covariance}. Risk set lúc train là minibatch vật lý; tích lũy gradient không ghép chúng thành risk set toàn cohort.

\begin{table}[htbp]
\centering
\small
\caption[Giao thức của các nhóm mô hình]{Các nhóm giao thức trong tiến trình phát triển mô hình. Cùng seed và split chưa có nghĩa là mọi cấu hình huấn luyện đều giống nhau.}
\label{tab:protocol-groups}
\begin{tabularx}{\linewidth}{@{}Y C{1.7cm} L{5.3cm}@{}}
\toprule
\textbf{Nhóm mô hình} & \textbf{Cox batch} & \textbf{Phạm vi so sánh} \\
\midrule
MRI Only; Longitudinal MRI & 6 & Khảo sát biểu diễn MRI và thời gian \\
Simple Multimodal Fusion & 8 & Mốc phát triển bằng ghép nối \\
Dynamic/MedicalNet Pairwise & 6 & So sánh initialization có kiểm soát \\
CP/RC-Free Pairwise & 6 & So sánh trong cùng họ pairwise \\
Ablation của CP Pairwise & 6 & Chỉ thay một thành phần so với CP đầy đủ \\
CNN--PCA--T-LSTM Adaptation & 16 & Pipeline tham chiếu, tối ưu riêng trên train \\
\bottomrule
\end{tabularx}
\end{table}

\section{MRI một thời điểm và MRI dọc theo thời gian}
\label{sec:mri-baselines}
\label{sec:longitudinal-results}
MRI Only dùng MRI tại visit cuối để dự báo nguy cơ. Trong so sánh 3D ResNet-18 cùng seed, Random ResNet đạt validation 0,787500, cao hơn MedicalNet 0,733333, nhưng test hậu nghiệm của MedicalNet lại cao hơn. Sự đảo thứ hạng này cho thấy không nên chọn initialization bằng cách nhìn lại test. CrossFormer3D, được thích nghi từ hướng CrossFormer \cite{wang2022crossformer}, giảm số tham số MRI Only từ 33.051.969 xuống 9.856.501 và đạt validation 0,764583; đây là một hướng khảo sát backbone gọn hơn.

\begin{table}[htbp]
\centering
\small
\setlength{\tabcolsep}{4pt}
\caption[C-index của các mô hình MRI trên cùng seed]{MRI một thời điểm và MRI dọc ở seed 20260727. Val/Test là C-index; Test là đánh giá hậu nghiệm. Dấu sao chỉ run attention còn hạn chế về hồ sơ hoàn tất.}
\label{tab:mri-milestones}
\begin{tabularx}{\linewidth}{@{}Y r r r@{}}
\toprule
\textbf{Mô hình} & \textbf{Epoch} & \textbf{Val $\uparrow$} & \textbf{Test $\uparrow$} \\
\midrule
MRI Only -- MedicalNet & 3 & 0,733333 & \bestresult{0,623106} \\
MRI Only -- Random ResNet & 19 & \bestresult{0,787500} & 0,496212 \\
MRI Only -- CrossFormer & 8 & 0,764583 & 0,490530 \\
Longitudinal MRI -- T-LSTM & 32 & 0,733333 & 0,492424 \\
Longitudinal MRI + Attention* & 34 & 0,783333 & 0,577652 \\
Longitudinal MRI -- CrossFormer & 25 & 0,725000 & 0,606061 \\
\bottomrule
\end{tabularx}
\end{table}

Longitudinal MRI dùng chung encoder cho ba visit, sau đó đưa chuỗi embedding qua T-LSTM. Phiên bản lấy trạng thái cuối đạt validation 0,733333; bổ sung attention ghi nhận 0,783333. Run attention được sửa sau lỗi AMP overflow, nhưng hồ sơ hoàn tất trong nguồn lịch sử còn hạn chế, nên được đánh dấu sao dù báo cáo đánh giá mới ghi nhận chạy thành công. Longitudinal MRI với CrossFormer đạt 0,725000. Các mốc này chưa chứng minh rằng thêm MRI theo thời gian tự động cải thiện dự báo; cần kiểm tra biểu diễn thời gian trên nhiều seed và điều kiện huấn luyện đồng nhất.

\section{So sánh các cách hợp nhất thông tin}
\label{sec:concat-results}
\label{sec:dynamic-pairwise-results}
Simple Multimodal Fusion ghép các embedding trước T-LSTM, attention và Cox head, đạt validation 0,812500 và test hậu nghiệm 0,723485. Dynamic Pairwise Fusion mô hình hóa MR, MC và RC bằng outer-product sau chiếu giảm chiều, rồi đưa tương tác trở lại từng nhánh qua gate, hệ số $\lambda$ và identity residual; validation tăng lên 0,828125. Tuy nhiên, Cox batch thay đổi từ 8 sang 6, nên mức tăng này là quan sát trong quá trình phát triển, chưa phải ablation chỉ thay cách fusion.

MedicalNet Pairwise Fusion giữ cấu trúc và giao thức của Dynamic Pairwise, chỉ đổi initialization MRI, đạt validation 0,839583. Đây là đối chứng full outer-product dùng cho thí nghiệm CP. Các biến thể CrossFormer được trình bày cùng tên của cơ chế hợp nhất trong Bảng~\ref{tab:fusion-common-seed}; kết quả test của chúng không được dùng để thay đổi lựa chọn kiến trúc sau khi đã chốt theo validation.

\begin{table}[htbp]
\centering
\small
\setlength{\tabcolsep}{4pt}
\caption[C-index của các cách fusion trên cùng seed]{Các cách fusion ở seed 20260727; Test là C-index hậu nghiệm của checkpoint chọn theo validation. Các nhóm khác giao thức không được coi là ablation chỉ thay fusion.}
\label{tab:fusion-common-seed}
\begin{tabularx}{\linewidth}{@{}Y r r r@{}}
\toprule
\textbf{Mô hình} & \textbf{Epoch} & \textbf{Val $\uparrow$} & \textbf{Test $\uparrow$} \\
\midrule
Simple Multimodal Fusion & 29 & 0,812500 & 0,723485 \\
CrossFormer Concatenation & 22 & 0,806250 & 0,876894 \\
Dynamic Pairwise Fusion & 5 & 0,828125 & 0,844697 \\
MedicalNet Pairwise Fusion & 35 & \bestresult{0,839583} & 0,763258 \\
CP Pairwise Fusion & 8 & 0,821875 & \bestresult{0,884470} \\
CrossFormer Static Pairwise & 9 & 0,804167 & 0,861742 \\
CrossFormer Dynamic Pairwise & 8 & 0,825000 & 0,878788 \\
CrossFormer Gated Residual Pairwise & 28 & 0,830208 & 0,867424 \\
\bottomrule
\end{tabularx}
\end{table}

\begin{figure}[htbp]
\centering
\includegraphics[width=\linewidth]{performance-evolution.pdf}
\caption[Validation theo các nhóm giao thức phát triển]{Các mốc validation ở seed 20260727, tách theo giao thức phát triển. Concatenation dùng Cox batch 8, các nhóm ResNet còn lại dùng batch 6; không nối các nhóm thành một xu hướng tăng điểm có kiểm soát. Dấu sao và marker rỗng chỉ run attention còn hạn chế về hồ sơ hoàn tất. Hình chỉ biểu diễn validation.}
\label{fig:performance-evolution}
\end{figure}

\FloatBarrier
\section{Đối chiếu với CNN--PCA--T-LSTM Adaptation}
\label{sec:paper-adaptation-results}
Kiến trúc tham chiếu của Aghajanian và cộng sự dùng CNN để trích embedding MRI dọc, PCA, T-LSTM và temporal attention cho dự báo sống còn \cite{aghajanian2025longitudinal}. Bản triển khai ở đây được gọi là \emph{CNN--PCA--T-LSTM Adaptation}, vì cohort, tiền xử lý và cấu hình huấn luyện không trùng hoàn toàn với bài báo. Cấu hình T-LSTM được chọn trong train, rồi khóa cho ba seed; validation và test bên ngoài không được dùng trong bước retuning đó.

CNN của từng seed được huấn luyện riêng và cố định để trích embedding 512 chiều trước dropout. PCA99 fit trên train tạo 38, 69 và 2 thành phần cho ba seed; T-LSTM dùng hidden size 64, Adam và batch 16 (Bảng~\ref{tab:alternative-encoder-shapes}). Ở seed 20260727, adaptation đạt validation 0,662500 và test hậu nghiệm 0,537879. Bảng~\ref{tab:baseline-gains} ghi mức tăng C-index so với bản adaptation này; ký hiệu $\uparrow$ thể hiện chiều tốt hơn của chỉ số, không biểu thị ý nghĩa thống kê.

\begin{table}[htbp]
\centering
\small
\setlength{\tabcolsep}{4pt}
\caption[C-index và mức tăng so với CNN--PCA--T-LSTM Adaptation]{C-index và mức tăng so với CNN--PCA--T-LSTM Adaptation trên cùng seed 20260727. Test là hậu nghiệm; mức tăng chỉ áp dụng cho bản adaptation trên cohort hiện tại.}
\label{tab:baseline-gains}
\begin{tabularx}{\linewidth}{@{}Y r r r r@{}}
\toprule
\textbf{Mô hình} & \textbf{Val $\uparrow$} & \textbf{Test $\uparrow$} & \textbf{$\Delta$ Val $\uparrow$} & \textbf{$\Delta$ Test $\uparrow$} \\
\midrule
CNN--PCA--T-LSTM Adaptation & 0,662500 & 0,537879 & $+0{,}000000$ & $+0{,}000000$ \\
MedicalNet Pairwise Fusion & 0,839583 & 0,763258 & $+0{,}177083$ & $+0{,}225379$ \\
CP Pairwise Fusion & 0,821875 & \bestresult{0,884470} & $+0{,}159375$ & $+0{,}346591$ \\
RC-Free Pairwise Fusion & \bestresult{0,841667} & 0,875000 & $+0{,}179167$ & $+0{,}337121$ \\
RC-Free Fixed-$\lambda$ & 0,833333 & 0,850379 & $+0{,}170833$ & $+0{,}312500$ \\
\bottomrule
\end{tabularx}
\end{table}

RC-Free tăng validation 0,179167 và test hậu nghiệm 0,337121 so với adaptation trên split hiện tại. Đây là so sánh giữa hai pipeline có đầu vào và cách huấn luyện khác nhau, nên không thể quy toàn bộ mức tăng cho riêng CP hoặc một route. Bảng cũng không đối chiếu trực tiếp với điểm công bố trên cohort của bài báo gốc.

\section{CP Pairwise Fusion}
\label{sec:cp-results}
CP Pairwise Fusion thay lớp song tuyến đầu tiên của full outer-product bằng $O(A^{\T}x\Had B^{\T}y)+b$. Encoder, post-MLP, directed projection, gate, $\lambda$, residual và mô hình thời gian được giữ nguyên; ranks MR/MC/RC là 8/8/4. CP tham số hóa tensor trọng số, thay vì phân rã dữ liệu của người bệnh.

Ở seed 20260727, CP đạt validation 0,821875 và test hậu nghiệm 0,884470; MedicalNet Pairwise đạt 0,839583 và 0,763258. CP có test cao hơn nhưng validation thấp hơn, nên không thể dùng cột test để tuyên bố CP là lựa chọn tốt hơn. Qua ba seed validation, full outer-product đạt $0{,}834375\pm0{,}014768$, còn CP đạt $0{,}827778\pm0{,}025007$; mean giảm 0,006597. Lợi ích đo được của CP là lõi tương tác ít tham số hơn, trong khi validation vẫn ở mức cạnh tranh.

\section{Ablation: vai trò của từng thành phần fusion}
\label{sec:ablation-results}
Để phân biệt tác động của từng thành phần, ablation lấy CP Pairwise Fusion đầy đủ làm reference và chỉ thay một yếu tố ở seed 20260727. Bỏ gate vẫn giữ $\lambda$ học được; cố định $\lambda=0{,}1$ vẫn giữ gate; bỏ identity residual chỉ bỏ đường truyền trực tiếp embedding gốc. Với biến thể bỏ MR, MC hoặc RC, toàn bộ pair block tương ứng và cả hai route có hướng của cặp đó đều bị bỏ, còn các encoder đầu vào vẫn được giữ.

Độ thay đổi validation so với CP đầy đủ là $\Delta_{\mathrm{Val}}$; chênh lệch giữa hai tập được tính bằng
\begin{equation}
\delta=C_{\mathrm{test}}-C_{\mathrm{validation}}.
\label{eq:split-gap}
\end{equation}
$\delta$ chỉ mô tả hai điểm số của một checkpoint trên hai tập dữ liệu. Giá trị gần 0 không tự chứng minh độ ổn định, bởi hai tập có nhóm người và cặp so sánh khác nhau.

\begin{table}[htbp]
\centering
\small
\setlength{\tabcolsep}{4pt}
\caption[Ablation của CP Pairwise Fusion]{Ablation cùng seed 20260727. Test là C-index hậu nghiệm; $\Delta$ Val so với CP đầy đủ, còn $\delta=\mathrm{Test}-\mathrm{Val}$. Không dùng $\delta$ để chọn mô hình.}
\label{tab:ablation}
\begin{tabularx}{\linewidth}{@{}Y r r r r@{}}
\toprule
\textbf{Biến thể} & \textbf{Val $\uparrow$} & \textbf{Test $\uparrow$} & \textbf{$\Delta$ Val} & \textbf{$\delta$} \\
\midrule
CP Pairwise Fusion & 0,821875 & \bestresult{0,884470} & $+0{,}000000$ & $+0{,}062595$ \\
CP -- Bỏ gate & 0,820833 & 0,867424 & $-0{,}001042$ & $+0{,}046591$ \\
CP -- Cố định $\lambda$ & 0,836458 & 0,862689 & $+0{,}014583$ & $+0{,}026231$ \\
CP -- Bỏ identity residual & 0,827083 & 0,776515 & $+0{,}005208$ & $-0{,}050568$ \\
CP -- Bỏ MR & 0,816667 & 0,842803 & $-0{,}005208$ & $+0{,}026136$ \\
CP -- Bỏ MC & 0,837500 & 0,857008 & $+0{,}015625$ & $+0{,}019508$ \\
RC-Free -- Bỏ RC & \bestresult{0,841667} & 0,875000 & $+0{,}019792$ & $+0{,}033333$ \\
\bottomrule
\end{tabularx}
\end{table}

Trong sáu biến thể thay đổi thành phần, RC-Free có cả validation và test hậu nghiệm cao nhất, lần lượt là 0,841667 và 0,875000. So với CP đầy đủ, validation tăng 0,019792, còn test giảm 0,009470; độ chênh tuyệt đối giữa test và validation giảm từ 0,062595 xuống 0,033333. Như vậy, RC-Free cho sự kết hợp thuận lợi giữa validation cao, test hậu nghiệm cao và cấu trúc ít cặp hơn, nhưng không đạt test cao nhất nếu tính cả reference CP đầy đủ.

\begin{figure}[H]
\centering
\includegraphics[width=\linewidth]{ablation-deltas.pdf}
\caption[Thay đổi validation trong ablation]{Thay đổi validation C-index so với CP Pairwise Fusion đầy đủ 0,821875 ở seed 20260727. Hình chỉ biểu diễn validation; không có error bar hay kiểm định ý nghĩa thống kê.}
\label{fig:ablations}
\end{figure}

Bỏ MR làm giảm validation, còn bỏ MC và RC làm tăng điểm ở seed này; bỏ identity residual cũng làm test giảm rõ xuống 0,776515. Mức test giảm khi bỏ residual là quan sát hậu nghiệm. Cấu trúc RC-Free được ưu tiên theo validation của ablation và mục tiêu rút gọn tương tác; các số test không được dùng để chọn lại thành phần. Các biến thể cố định $\lambda$, bỏ MR và bỏ MC còn có $|\delta|$ nhỏ hơn RC-Free, vì vậy không dùng riêng chênh lệch test--validation để xếp hạng độ ổn định.

\FloatBarrier
\section{RC-Free Pairwise Fusion và kiểm tra nhiều seed}
\label{sec:rcfree-multiseed}
RC-Free Pairwise Fusion giữ đủ ba nhánh đầu vào, chỉ bỏ cặp RC và hai route $(\mathrm{RC}\to R)$, $(\mathrm{RC}\to C)$. Hai cặp MR/MC, bốn dynamic gate, bốn hệ số $\lambda$ học được, identity residual, T-LSTM và attention được giữ lại. Kiến trúc được chọn theo validation của ablation; test hậu nghiệm dùng để mô tả checkpoint đã chọn.

Validation trung bình ba seed đạt 0,834722 với sample SD 0,023693, cao hơn CP đầy đủ 0,006944 và gần MedicalNet Pairwise, với chênh lệch mean chỉ 0,000347. RC-Free không tốt hơn ở mọi seed: seed 20260729 đạt 0,808333, thấp hơn MedicalNet Pairwise 0,817708. Bằng chứng nhiều seed vì vậy hỗ trợ một cấu trúc gọn có hiệu năng cạnh tranh, thay vì chứng minh ưu thế tuyệt đối.

RC-Free Fixed-$\lambda$ là phép kiểm tra tiếp theo trên cấu trúc chỉ giữ MR/MC; nó khác biến thể cố định $\lambda$ của CP đầy đủ trong bảng ablation. Ở seed 20260727, mô hình này chọn epoch 12, đạt validation 0,833333 và test hậu nghiệm 0,850379. Mean validation ba seed bằng RC-Free, nhưng SD là 0,031273. Việc cố định $\lambda$ chưa tạo ra lợi ích rõ ràng, nên kiến trúc cuối giữ các hệ số học được.

\section{Benchmark validation chuẩn hóa trên ba seed}
\label{sec:standardized-benchmark}
Benchmark đối chiếu các best checkpoint theo ba seed 20260727, 20260728 và 20260729, kiểm tra fingerprint của split, preprocessing train-only và quy ước C-index. Cả 15 ô model--seed được replay từ các run đã tồn tại. Pipeline adaptation vẫn giữ cách huấn luyện riêng; việc chuẩn hóa đánh giá không biến tất cả model family thành một ablation duy nhất.

\begin{table}[htbp]
\centering
\small
\setlength{\tabcolsep}{4pt}
\caption[Benchmark validation chuẩn hóa trên ba seed]{C-index validation chuẩn hóa: mean $\pm$ sample SD ($n=3$). Hai dòng in đậm đồng hạng theo mean. Seed 27/28/29 là 20260727/20260728/20260729.}
\label{tab:standard-bench}
\begin{tabularx}{\linewidth}{@{}Y r r r r@{}}
\toprule
\textbf{Mô hình; C-index $\uparrow$} & \textbf{Seed 27} & \textbf{Seed 28} & \textbf{Seed 29} & \textbf{Mean $\pm$ SD} \\
\midrule
Paper CNN--T-LSTM & 0,662500 & 0,770833 & 0,791667 & $0{,}741667\pm0{,}069347$ \\
MedicalNet Pairwise & 0,839583 & 0,845833 & 0,817708 & $0{,}834375\pm0{,}014768$ \\
CP Pairwise & 0,821875 & 0,855208 & 0,806250 & $0{,}827778\pm0{,}025007$ \\
\bestresult{RC-Free Pairwise} & \bestresult{0,841667} & \bestresult{0,854167} & \bestresult{0,808333} & \bestresult{$0{,}834722\pm0{,}023693$} \\
\bestresult{RC-Free Fixed-$\lambda$} & \bestresult{0,833333} & \bestresult{0,866667} & \bestresult{0,804167} & \bestresult{$0{,}834722\pm0{,}031273$} \\
\bottomrule
\end{tabularx}
\end{table}

Paper CNN--T-LSTM trong bảng là CNN--PCA--T-LSTM Adaptation; MedicalNet Pairwise là full outer-product, CP Pairwise giữ đủ ba cặp, còn RC-Free Pairwise chỉ giữ MR/MC. Hai dòng in đậm đồng hạng theo mean; sample SD dùng mẫu số $n-1$. RC-Free cao hơn adaptation về mean 0,093056, nhưng ba seed dùng cùng split, không phải ba cohort độc lập.

\begin{figure}[H]
\centering
\includegraphics[width=\linewidth]{standardized-benchmark.pdf}
\caption[Benchmark validation chuẩn hóa trên ba seed]{Mean và sample SD validation của năm mô hình. Error bar biểu diễn dao động giữa các seed trên cùng split, không phải khoảng tin cậy hay kết quả test.}
\label{fig:benchmark}
\end{figure}

MedicalNet Pairwise có SD thấp nhất trong nhóm pairwise, còn RC-Free và RC-Free Fixed-$\lambda$ đồng hạng mean. Với chỉ ba seed và nhiều lần dùng validation để chọn kiến trúc, chưa đủ bằng chứng kết luận RC-Free ổn định nhất. Khóa luận chọn RC-Free là cấu hình phù hợp nhất với tiêu chí validation và mục tiêu giữ cấu trúc gọn trong nhóm đã khảo sát.

\FloatBarrier
\section{Hiệu quả về số tham số}
\label{sec:computational-efficiency}
Ba phép biến đổi song tuyến đầu tiên của full outer-product có 229.696 tham số gồm bias. CP giảm chúng còn 3.712, tương đương 98,384\%. Các post-MLP không thay đổi; tính cả post-MLP, tổng của ba pair block giảm từ 248.608 còn 22.624. Bảng~\ref{tab:interaction-parameters} phân biệt lõi interaction với toàn bộ mạng để tránh diễn giải mức giảm 98\% như giảm toàn mô hình.

\begin{table}[htbp]
\centering
\small
\caption[Số tham số của full outer-product và CP]{Số tham số phép biến đổi song tuyến đầu tiên, gồm bias. CP ranks MR/MC/RC là 8/8/4.}
\label{tab:interaction-parameters}
\begin{tabularx}{\linewidth}{@{}Y r r@{}}
\toprule
\textbf{Phạm vi} & \textbf{Full outer-product} & \textbf{CP} \\
\midrule
MR: $32\times32\to128$ & 131.200 & 1.664 \\
MC: $32\times16\to128$ & 65.664 & 1.536 \\
RC: $32\times16\to64$ & 32.832 & 512 \\
\midrule
Tổng ba lớp đầu & 229.696 & 3.712 \\
Ba pair block, gồm post-MLP & 248.608 & 22.624 \\
Toàn mô hình & 33.760.941 & 33.534.957 \\
\bottomrule
\end{tabularx}
\end{table}

\begin{figure}[H]
\centering
\includegraphics[width=\linewidth]{cp-parameter-reduction.pdf}
\caption[Tham số outer-product và CP theo từng phạm vi]{Giảm tham số từ full outer-product sang CP: 98,384\% ở các lớp song tuyến đầu tiên, nhưng chỉ 0,669\% ở toàn mô hình do MRI encoder chiếm phần lớn tham số. Post-MLP được giữ nguyên.}
\label{fig:cp-parameters}
\end{figure}

RC-Free có tổng 33.519.497 tham số, trong đó hai pair block MR/MC có 19.968 tham số, lõi CP có 3.200, bốn gate có 25.348 và bốn $\lambda$ là bốn scalar học được. Fixed-$\lambda$ còn 33.519.493 tham số. Việc bỏ RC chủ yếu đơn giản hóa interaction graph; mức giảm toàn mạng là nhỏ. Parameter count không thay thế phép đo FLOPs, latency hoặc memory. Các run có số epoch dừng khác nhau nên thời gian train tổng không được dùng như một benchmark tốc độ có kiểm soát.

\FloatBarrier

\section{Kết luận từ thực nghiệm và lựa chọn kiến trúc}
\label{sec:rq-results}
\begin{table}[H]
\centering
\small
\caption[Kết luận cho các câu hỏi nghiên cứu]{Trả lời các câu hỏi nghiên cứu trong phạm vi dữ liệu và giao thức hiện có.}
\label{tab:rq-answer}
\begin{tabularx}{\linewidth}{@{}C{0.9cm}Y@{}}
\toprule
\textbf{RQ} & \textbf{Kết luận và giới hạn} \\
\midrule
RQ1 & Các mốc longitudinal MRI chưa cho thấy lợi ích ổn định so với MRI Only; cần so sánh biểu diễn thời gian trong điều kiện đồng nhất. \\
RQ2 & Simple Multimodal Fusion có validation cao hơn longitudinal MRI, nhưng khác batch và hồ sơ attention hạn chế chưa cho phép tách riêng đóng góp của đầu vào. \\
RQ3 & Pairwise có validation cao hơn concatenation trong tiến trình phát triển; cần đối chứng cùng batch để cô lập tác động của fusion. \\
RQ4 & CP giảm 98,384\% tham số ở các lớp song tuyến đầu tiên; mean validation giảm 0,006597 so với full outer-product. \\
RQ5 & RC-Free có validation cao nhất trong ablation cùng seed, test hậu nghiệm cao và mean validation cạnh tranh qua ba seed; giữ MR/MC là lựa chọn thiết kế được hỗ trợ trong cohort này. \\
\bottomrule
\end{tabularx}
\end{table}

Kiến trúc chính của khóa luận là \textbf{RC-Free Pairwise Fusion}. Quyết định dựa trên validation và mục tiêu đơn giản hóa tương tác; mức test hậu nghiệm 0,875000 bổ sung bằng chứng mô tả, không thay đổi tiêu chí lựa chọn. CP đầy đủ có test cao hơn và MedicalNet Pairwise có SD validation thấp hơn, vì vậy kết luận không mở rộng thành ưu thế trên mọi chỉ số. Đánh giá trên cohort độc lập hoặc một tập giữ lại mới chưa được thực hiện và là bước cần thiết để kiểm chứng khả năng khái quát hóa.
